% Generated draft: requires mathematical and source review.
\documentclass[11pt]{article}
\usepackage[margin=25mm]{geometry}
\usepackage{fontspec}
\setmainfont{Latin Modern Roman}
\usepackage{amsmath,amssymb,mathtools,unicode-math}
\setmathfont{Latin Modern Math}
\usepackage{enumitem,needspace,xurl,hyperref}
\setlength{\parindent}{0pt}
\setlength{\parskip}{5pt}
\setlength{\emergencystretch}{3em}
\newcommand{\N}{\mathbb{N}}
\newcommand{\Z}{\mathbb{Z}}
\newcommand{\Q}{\mathbb{Q}}
\newcommand{\R}{\mathbb{R}}
\newcommand{\C}{\mathbb{C}}
\newcommand{\F}{\mathbb{F}}
\newcommand{\A}{\mathbb{A}}
\newcommand{\PP}{\mathbb{P}}
\newcommand{\E}{\mathbb{E}}
\newcommand{\eps}{\varepsilon}
\newcommand{\dd}{\,\mathrm{d}}
\begin{document}
\section*{UnsolvedMath ID: 700511}
\textbf{Source problem code:} catalogue entry
\subsection{Definitions and mathematical statement}
For a topological group \(G\), let
\[
\mathcal N_{\mathrm{cl}}(G)=\{N\leq G:N\text{ is closed in }G\text{ and }gNg^{-1}=N\text{ for every }g\in G\}.
\]
Write \(\tau_E(G)\) for the \(E\)-topology on this set referred to in the source. The supplied record does not define this topology.

The problem, proposed by I. V. Protasov, is to characterize all locally compact locally pronilpotent groups \(G\) for which
\[
\bigl(\mathcal N_{\mathrm{cl}}(G),\tau_E(G)\bigr)
\quad\text{is compact}.
\]
Explicitly, the compactness condition means that for every index set \(I\) and every family \((U_i)_{i\in I}\) of \(\tau_E(G)\)-open subsets satisfying
\[
\mathcal N_{\mathrm{cl}}(G)=\bigcup_{i\in I}U_i,
\]
there exists a finite subset \(J\subseteq I\) such that
\[
\mathcal N_{\mathrm{cl}}(G)=\bigcup_{i\in J}U_i.
\]
The term \emph{locally pronilpotent} is retained from the source: its precise local condition and the convention for pronilpotence are not specified in the supplied record. The record also does not explicitly specify a Hausdorff convention for locally compact groups. These conventions, together with the definition of \(\tau_E(G)\), must be recovered from the cited source before this becomes a fully specified classification problem.

The source states that the corresponding problem for discrete groups has been solved, but supplies neither that classification nor a reference for it.
\subsection{Short English statement}
Characterize the locally compact locally pronilpotent groups whose closed normal subgroups form a compact space in the source's \(E\)-topology. The supplied record leaves that topology and the precise meaning of local pronilpotence undefined. It reports that the discrete case has been solved.
\subsection{Sources}
\begin{itemize}
\item UnsolvedMath local registry, permanent ID \(700511\), exact source problem code \texttt{catalogue entry}; supplied upstream revision \texttt{d371ee917ac9647fbb06cd6fc529deec6f3dda82}.
\item \emph{The Kourovka Notebook (Unsolved Problems in Group Theory)}, 8th issue (1982), Problem 8.62; proposed by I. V. Protasov. The supplied record also identifies the 21st issue edition (2026), \texttt{arXiv:1401.0300v48}. Its collection description attributes editing to E. I. Khukhro and V. D. Mazurov.
\item The registry labels this problem open and reports that the corresponding discrete-group problem has been solved. These are supplied source assessments, not independently verified mathematical or literature conclusions. No reference for the discrete classification is supplied.
\end{itemize}
\end{document}
